% 阶段研究汇总（第二阶段）；AI 起草，待用户过目后定稿编译。
\documentclass[UTF8,11pt,a4paper]{ctexart}
\usepackage[margin=2.6cm]{geometry}
\usepackage{amsmath,amssymb}
\usepackage{booktabs}
\usepackage[colorlinks=true,linkcolor=blue,citecolor=blue,urlcolor=blue]{hyperref}
\usepackage[backend=biber,style=numeric,sorting=none,doi=true,url=false,eprint=true]{biblatex}
\renewcommand*{\bibfont}{\small}
\setlength{\bibitemsep}{0pt}
\addbibresource{../../literature/references.bib}

\title{第二阶段：有限温三临界结构与全局证据\\[0.4em]
\large 热三临界候选、相图拓扑与证据分级}
\author{Quantum Optics Thesis 工作笔记}
\date{2026 年 11 月 26 日}

\begin{document}
\maketitle

\begin{abstract}
这一阶段做的是：在 $zz$ 交换、$z$ 纵场、$x$ 腔耦合的量子 Dicke--Ising 链里，看零温的
一级/二级边界和三临界点怎么延伸到有限温。主要结果有三个：$J=-0.2,\Delta=0.2$ 处找到
热三临界候选 $(T^\ast,g^\ast)=(0.2572636,0.4888020)$，六阶系数恒正、小场按 $h^4$ 标度，
属于数值强支持；多 $\Delta$ 的高温根排成非单调上支，低温在 $0.4<\Delta<0.4002585$ 出现
由上下双根夹住的重入一级窄袋，两根在 $(0.4002585,0.1255994)$ 切触；其中 $\Delta=0.4001$
一段截面拿到全局数值支持（分辨率 68.85），而切触端点仍只是局部候选。全部来自平衡有效
哈密顿量，不涉及开放系统。
\end{abstract}

\section{背景与问题}

光与物质集体耦合达到临界强度会发生超辐射相变，但能不能在真实腔里实现受求和规则约束
\cite{Nataf2010NoGo}；开放腔里已经看到 Dicke 型转变\cite{Baumann2010OpenDicke}。这类模型
零温和有限温的一般理论早有人系统梳理\cite{Peng2019Unified}。加上 Ising 交换后，磁序和
光场序可以共存，电路 QED 框架下就预测过正常、超辐射、反铁磁正常、反铁磁超辐射共存这些相
\cite{Zhang2014CircuitIsing}。

零温这一侧，近几年已经被做得很透：Langheld 等用大尺度 QMC 给出定量相图、认出狭窄的反铁磁--
超辐射共存区\cite{Langheld2025Wormhole}，Leibig 等用自洽纯自旋模型加 NLCE/DMRG 独立确认
\cite{Leibig2026SelfConsistent}，Koziol 给出铁磁三临界点的 Landau 展开\cite{Koziol2026Criticality}，
Hörmann 等进一步报零温解析相界和三临界条件\cite{Hormann2026Folds}。围绕 Mendonça 等的
Ising/XXZ 基态结论还有一轮 Comment/Reply 争议\cite{Mendonca2025Matter,Hormann2025Comment}。
所以零温相图本身不剩什么新东西，重画一遍不是创新。

有限温这一侧，已有的先例要么轴向不同、要么模型不同：Gammelmark--Mølmer 最早做量子 Ising 链
与单模腔的有限温自由能和一、二阶超辐射边界，但轴向不一样\cite{Gammelmark2011Thermal}；
Otake--Bamba 解的是同轴\emph{经典} Ising 链的有限温\cite{Otake2026FiniteT}；Rao 等做可积
$zz/y/x$ 链的有限温三临界\cite{Rao2025Unilateral}；Granet 做 $zz/x/x$ 的热转变\cite{Granet2023Exact}；
Román--Roche 等做同轴协同的有限温响应\cite{RomanRoche2021Photon}；相互作用多比特腔的一级跃迁
也有早先例\cite{Lee2004FirstOrder}。定向检索下来，非零纵场 $zz/z/x$ 这个同模型的完整热相图
还没人做过（索引漏收和 Leibig 期刊正文没核到，这点保留）。

所以这一阶段的问题很清楚：零温三临界点 $(\Delta,J)$ 附近，一级/二级边界往有限温延伸的时候，
是形成一个热三临界线和重入结构，还是只是边界平凡地平移一下？如果是后者，这个候选就该降级。

\section{模型与判据}

取 $\hbar=k_B=1$、$s_i^\alpha=\sigma_i^\alpha/2$、周期链、$J<0$ 铁磁，物质自旋链是
\begin{equation}
  H_{\rm spin}(h)=J\sum_i\sigma_i^z\sigma_{i+1}^z+\frac{\Delta}{2}\sum_i\sigma_i^z
  +\frac{h}{2}\sum_i\sigma_i^x,\qquad h=4g\varphi,\ \varphi=\mathrm{Re}\langle a\rangle/\sqrt N,
  \label{eq:Hspin}
\end{equation}
单模集体耦合做热力学自洽后，得到腔位移泛函
\begin{equation}
  \Phi_T(\varphi)=\omega\varphi^2+f_{\rm spin}(T,4g\varphi),\qquad
  f_{\rm spin}(T,h)=-\lim_{N\to\infty}\frac{T}{N}\log\operatorname{Tr}e^{-H_{\rm spin}(h)/T}.
  \label{eq:Phi}
\end{equation}
注意有限 $N$ 的完整光子--自旋配分函数不能直接拿 $\Phi_T$ 代替；一级边界必须比较
$\varphi\in[0,g/\omega]$ 里全部极小值的自由能。

正常态和某个有限位移态等能所需的耦合，做个变量替换就是
\begin{equation}
  g_{\rm equal}^2(h)=\frac{\omega h^2}{16\big[f_{\rm spin}(T,0)-f_{\rm spin}(T,h)\big]},
  \qquad g_{\rm onset}(T)=\inf_{h>0}g_{\rm equal}(T,h).
  \label{eq:gequal}
\end{equation}
配合 $0\le f_{\rm spin}(T,0)-f_{\rm spin}(T,h)\le h/2$，能证明凡是 $h>8g_{\rm ref}^2/\omega$
都不可能给出更小的等能耦合——这就给连续搜索一个有限的上界，不用扫到无穷。下确界在有限
$h$ 取得，就是一级跃迁候选；只在 $h\to0$ 取得，就退化成正常态局部失稳阈值 $g_{\rm local}$，
后者由 $h=0$ 的经典链转移矩阵加三自旋概率精确算出来。

小横场下 $f=f_0+a_2h^2+a_4h^4+a_6h^6+\cdots$，这些 Landau 系数用零场回返路径的合流差商
精确枚举。这里有个踩过的坑：最近邻双翻的中间能量差是
\begin{equation}
  dE_{ij}=dE_i+dE_j+4J s_i s_j,
  \label{eq:dEij}
\end{equation}
这一项从 $O(J^2)$ 起贡献，早期实现漏了它，导致 $T^\ast\approx0.3094$ 的旧结论算错，后来
撤回了（推导细节见 \texttt{manuscript/theory/model.md}）。

数值上用了四套互相独立的方法交叉：周期热 ED（$N\le12$，外部 32\,GB 机器扩到 $N=14$）、
开边界 NLCE（$L\le12$ 加 Aitken）、Suzuki--Trotter 无限空间转移矩阵（$M\le20$ 加 $1/M^2$
双窗口外推）、以及 60 位精度的精确 Dyson 回返路径。一级边界的 $g_{\rm equal}(h)$ 在有限域里
做粗扫、检出全部内点极小、黄金分割细化，再对每个极小做三方法诊断（细节和误差边界见
\texttt{manuscript/theory/numerics.md}）。

\section{热三临界候选}

先做零温校准。$\Delta=0$ 是横场 Ising 可积极限，全局一级点 $g_c=0.4090436682$，和 Koziol 的
数值逐位一致\cite{Koziol2026Criticality}。$\Delta>0$ 时小场展开给 $g_{\rm local}^2=\omega(\Delta-4J)/4$，
四次系数在 $\Delta=-2J$ 变号，零温三临界点 $g_{tc}^2=-3\omega J/2$，$J=-0.2$ 时就是
$(0.4,\sqrt{0.3})$——和 Langheld 等的相图对得上\cite{Langheld2025Wormhole}。零温一级线
$g_c(\Delta)$ 在 $\Delta=0.15$--$0.4$ 单调升到 $\sqrt{0.3}$、深度单调降到零。

然后是核心。对 $J=-0.2,\Delta=0.2,\omega=1$，修正后的无限链四次根给
\begin{equation}
  T^\ast=0.2572636163,\qquad g^\ast=g_{\rm local}(T^\ast)=0.4888020431,
  \label{eq:candidate}
\end{equation}
$a_2(T^\ast)=-0.2615856960$，$a_4(T^\ast)=0$（二分残差括区宽 $7.3\times10^{-14}$）。精确
有限环回返路径给 $a_6(N{=}6,8,10,12)=0.0371731,0.0374355,0.0374869,0.0374967$，末两尺寸差
$9.8\times10^{-6}$，恒正——这正是热三临界点局部存在的必要条件（二次、四次系数同时为零、
六次为正）。

候选点附近小场应该按
\begin{equation}
  g_{\rm equal}(T^\ast,h)-g_{\rm local}
  =\frac{g_{\rm local}\,a_6}{2(-a_2)}\,h^4+O(h^6),
  \label{eq:h4}
\end{equation}
走，$N=12$ 预测系数 $0.0350334$。无限空间转移矩阵在 $h=0.025,0.05,0.1$ 给的差
$1.3\times10^{-8},2.18\times10^{-7},3.49\times10^{-6}$ 正好对上这个 $h^4$；$T=0.25<T^\ast$
小场差为负、$T=0.265>T^\ast$ 为正，两侧符号切换被周期 ED 和无限空间转移矩阵一起确认。
有限域 $0<h\le8g_{\rm local}^2/\omega\approx1.911$ 的全域驻点扫描里，$T=0.245,0.25,0.255$
的内点极小深度按 Landau 预测
\begin{equation}
  \mathrm{depth}\approx g_{\rm local}\,\frac{a_4^2}{8a_6\,|a_2|},\qquad h^{\ast\,2}=-\frac{a_4}{2a_6}
  \label{eq:depth}
\end{equation}
塌缩（实测 $1.23\times10^{-4},4.39\times10^{-5},4.05\times10^{-6}$，比值 $1.01/1.00/0.93$），
而 $T\ge T^\ast$ 没有低于 $g_{\rm local}$ 的采样点。六阶收敛、$h^4$ 标度、邻域符号切换、有限域
驻点四条合起来，这个点可以叫\emph{数值强支持的热三临界候选}——但离散扫描排除不了任意窄的
远场极小，所以还不能叫已证明。

\section{相图拓扑}

对多个 $\Delta$ 求 $a_4(T)=0$ 的根，高温上支 $T_{\rm upper}(\Delta)$ 排成一个非单调的拱：
$\Delta=0.15,0.2,0.25,0.3,0.35,0.38,0.4$ 对应
$T_{\rm upper}=0.252465,0.257264,0.255727,0.247082,0.227322,0.202548,0.141122$。
在 $\Delta=0.15,0.2$（完整三方法）和 $0.3,0.38$（轻量）四个切片上，上支下侧（一级侧）
都有低于 $g_{\rm local}$ 的内点极小、深度按式~\eqref{eq:depth} 塌缩，上侧（连续侧）没有
内点极小——对上"一级→连续"这个结构。

低温这边更细：全区间多根搜索发现 $0.4<\Delta<0.4002585$ 里 $a_4(T)$ 按正→负→正变化，
上下两根夹出一个重入一级窄袋，两根随 $\Delta$ 增大越靠越近，最后在
\begin{equation}
  (\Delta_{\rm end},T_{\rm end},g_{\rm end})=(0.4002585148,0.1255994291,0.5477461329)
  \label{eq:endpoint}
\end{equation}
切触。端点由无限链和有限环 $N=8,10,12$ 独立收敛，端点
$a_6=0.00886088613,0.00886088672,0.00886088672>0$。整个拓扑是：零温三临界点 $(\Delta=0.4,T=0)$
锚住低温下支；$\Delta<0.4$ 的一级区从 $T=0$ 一直延伸到上支；$0.4<\Delta<\Delta_{\rm end}$ 只在
上下两根之间有一级窄袋（重入）。

\section{证据分级}

重入窄袋的深度随 $\Delta\to\Delta_{\rm end}$ 掉得极快，所以必须按"方法差对信号的分辨率"逐点
分级，不能把局部候选直接写成全局相界：

\begin{center}
\begin{tabular}{lccc}
\toprule
$\Delta$ & 六阶预测深度 & ED/NLCE 实测 & 信号/经验分辨率 \\
\midrule
$0.4001$ & $1.92\times10^{-9}$ & $-1.898/-1.864\times10^{-9}$ & $57.3\to68.85$ \\
$0.4002$ & $2.62\times10^{-10}$ & 同号 & $1.42\to1.32$ \\
$0.40024$ & $2.62\times10^{-11}$ & $N=12\to14$ 翻号 & $<1$ \\
\bottomrule
\end{tabular}
\end{center}

$\Delta=0.4001$ 的 $N=14$ 周期 ED（外部 32\,GB 机器，源码和输入哈希全对上）给深度
$1.891849\times10^{-9}$，是六阶预测的 $98.3\%$，而且 $N=8,10$ 从 $h=0.03$ 扫到严格上界
约 $2.40$ 都没有远场内点极小——这一截面可以叫\emph{全局数值支持}。$\Delta=0.4002$ 还只是
边缘，$\Delta=0.40024$ 在 $N=12\to14$ 翻号，转移矩阵外推在 $10^{-11}$ 深度下的窗口漂移约
$3.9\times10^{-6}$，大出 3--5 个量级，没有裁决力。所以切触端点式~\eqref{eq:endpoint} 仍只是
\emph{局部 Landau 候选}。

\section{还差什么与下一步}

几个要说清楚的边界。第一，所有近等能误差都是\emph{数值诊断级}（ED 尺寸收敛、NLCE 与 ED 差、
转移矩阵两窗口差），不是严格误差界。第二，固定 $\varphi$ 的一维短程链在 $T>0$ 没有自身有限温
相变\cite{Araki1969Gibbs}，所以零温的交错磁长程序不能当 $T>0$ 的独立共存热相；但完整腔系统的
超辐射热转变不受这个限制。第三，零温的三临界点、窄共存区这些跟 Langheld、Leibig、Koziol、
Hörmann 一致\cite{Langheld2025Wormhole,Leibig2026SelfConsistent,Koziol2026Criticality,Hormann2026Folds}，
不是新东西；本工作的候选增量只在有限温这一侧的拓扑，而且"空白"本身还是可反证的假设。

下一步三件事：补核 Leibig 期刊正文和最新引文；真想把切触端点写成强主张，先立一个能处理
$10^{-10}$--$10^{-11}$ 深度的高精度响应/有限温 MPS/SSE 基准，而不是继续堆同类双精度 ED 的
尺寸；如果独立方法不收敛或撞见同模型完整先例，按主张台账 C04 降级或止损。

\printbibliography[title={参考文献}]
\end{document}
